\begin{proof}
\begin{enumerate}[label=\arabic*.]
\item Reduction modulo exact differentials is completely explicit.  To
differentiate $\omega_k$, solve the polynomial identity
\[
-\tfrac12x^k\partial_\lambda f
 =f\sum_{j=0}^3G_{kj}x^j+R_k'f-\tfrac12R_kf',
 \qquad\deg R_k\leq4,
\]
where $f=x^5+\lambda^3x+\lambda^4$.  The last two terms are the
numerator of $d(R_k/y)$.  Matching the nine coefficients gives the
four-by-four Gauss--Manin matrix exactly.
At the branch point at infinity, a local coordinate $z=x^{-1/2}$
gives $\omega_k=-2z^{2-2k}(1+O(z^8))\,dz$.
Residues of a primitive times a differential give the alternating
pairings $J_{03}=4/3$, $J_{12}=4$, with their negatives in reverse
order and all other entries zero.  This form is nondegenerate.
The invariant bivector is therefore $3\omega_{03}+\omega_{12}$.
Taking the exterior-square connection and this quotient gives
(1.15.2), as checked symbolically in \texttt{de\_rham\_connection.py}.

\item The root-chart powers computed from $x^kdx/y$ are
$u^{6k-9}$ at zero and $u^{6-4k}$ at infinity.
Their nondegenerate pairing shows that the four resulting regular
frames form the full de Rham bundle on those charts.  At the node
the computed connection has a simple pole and square-zero residue;
hence these are the nilpotent logarithmic extension frames there.
Thus the four line bundles are
$T^3,T,T^{-1},T^{-3}$, and their quotient exterior-square lines are
the stated powers of $L$.

\item All $H^1(L^j)$ for $j=0,\pm1,\pm2$ vanish by the coarse-degree
calculation already used in P1.8.  Consequently every holomorphic
extension of this $W$ by $\mathcal O$ splits as a bundle.
In such a splitting its flat connection is an upper block extension
by a row $b$.  Regularity at the finite root charts and a logarithmic
pole at the cusp require, entry by entry,
\[
b=\frac1{\lambda-\lambda_*}
 (a\lambda^{-2},b'\lambda^{-1},0,c,d\lambda).
\]
For example the $L^j$ entry, after $\lambda=u^4$, has exponent
$3j+3+4p$ for a term $\lambda^p/ (\lambda-\lambda_*)$;
the infinity chart requires $-4j-1-5p\geq0$.
These two inequalities leave exactly the four displayed terms.
The equation that the full cusp residue have square zero is
$b_{\rm res}(\Gamma_W)_{\rm res}=0$.
Its solution space has basis
\[
(-25/6,-125/48,1,0),\qquad(5/3,-25/3,0,1).
\]
The sole nonconstant holomorphic splitting change is the row selecting
$w_0$, and its connection coboundary is the third row of $\Gamma_W$.
In these two basis vectors that row is $(-48b_1+60b_2)/3125$.
Thus there is one nontrivial flat extension direction modulo splitting
changes and rescaling of the invariant line.  The first vector has
rank-three square-zero cusp residue, so it represents that direction.
The supplied extension is nontrivial, since its dual has no common
invariant vector by P1.1.  It is therefore isomorphic to (1.15.1).
This proves the connection identification without inferring it from
the subsequent numerical matching.

\item For the last normalization, the kernel of the quotient cusp residue
on $(w_2,w_1,w_0)$ is generated by
$v=(-4096/1171875,64/1875,1)^t$.
It satisfies $b_{\rm res}(v)=625/192$ and
$\alpha_{\rm alg}(v)=4/3$.
The logarithmic monodromy is $2\pi i$ times this residue.
The exact cusp identity of P1.8 characterizes the contraction on this
line by $N_V(\widetilde v)=\alpha(v)e$.
Since the space of global contraction forms is one-dimensional, the
ratio $(625/192)/(4/3)=625/256$ fixes its scalar everywhere and gives
the displayed $625\pi i/128$.

\end{enumerate}
\end{proof}